\documentclass[../main.tex]{subfiles}
\graphicspath{{\subfix{../images/}}}

\begin{document}
    \section{Conclusión}
    Durante el desarrollo de este proyecto se logró implementar un simulador
    modular de Máquinas de Turing capaz de validar cadenas pertenecientes a
    distintos lenguajes formales mediante una interfaz interactiva de línea de
    comandos. A través de este trabajo fue posible aplicar conocimientos
    fundamentales relacionados con teoría de la computación, lenguajes formales,
    estructuras de control y programación modular en lenguaje C.

    El proyecto permitió comprender de manera práctica el funcionamiento interno de
    una Máquina de Turing, especialmente en aspectos como el manejo de estados,
    transiciones, validación de símbolos y aceptación o rechazo de cadenas. Además,
    la división del programa en módulos independientes facilitó la organización del
    código y mejoró la claridad, mantenimiento y escalabilidad del sistema,
    permitiendo que cada archivo cumpliera una función específica dentro de la aplicación.

    Asimismo, la elaboración de diagramas, configuraciones y manuales ayudó a
    documentar correctamente el funcionamiento del programa y a comprender mejor la
    lógica utilizada durante el diseño e implementación. La creación de una
    interfaz basada en comandos también permitió simular de forma sencilla la
    interacción entre el usuario y las máquinas implementadas, facilitando la
    realización de pruebas para distintos lenguajes.

    En general, este proyecto no solo fortaleció habilidades técnicas de
    programación y análisis lógico, sino que también permitió relacionar conceptos
    teóricos vistos en clase con una aplicación funcional y práctica, demostrando
    la importancia de las Máquinas de Turing como base de la computación y del
    estudio de los lenguajes formales
    \end{document}
